\documentclass[aspectratio=169]{beamer}
\usepackage{tikz}
\usepackage{graphicx}
\usepackage{listings}
\usepackage{xcolor}
\usepackage{hyperref}
\usepackage[spanish,es-noquoting]{babel}
\usepackage[round]{natbib}

\lstset{
  basicstyle=\ttfamily\small,
  keywordstyle=\color{blue},
  stringstyle=\color{red},
  commentstyle=\color{gray},
  showstringspaces=false
}
% Minimal look
\usetheme{Warsaw}
%\setbeamertemplate{headline}
\setbeamertemplate{footline}{
  \leavevmode%
  \hbox{%
    % Left: Section name
    \begin{beamercolorbox}[wd=.33\paperwidth,ht=2.25ex,dp=1ex,left]{author in head/foot}%
      \hspace*{2ex}\insertsection
    \end{beamercolorbox}%
    % Center: Presentation title
    \begin{beamercolorbox}[wd=.34\paperwidth,ht=2.25ex,dp=1ex,center]{title in head/foot}%
      \insertshorttitle
    \end{beamercolorbox}%
    % Right: Slide numbers
    \begin{beamercolorbox}[wd=.33\paperwidth,ht=2.25ex,dp=1ex,right]{date in head/foot}%
      \insertframenumber{} / \inserttotalframenumber\hspace*{2ex}
    \end{beamercolorbox}%
  }%
  \vskip0pt%
}

\usecolortheme{default}
\usefonttheme{serif}
\graphicspath{{../}{./}}
\setbeamertemplate{caption}[numbered]

% Title info
\title[Módulo 6: El Límite es tu Imag]{Módulo 6: El Límite es tu Imaginación}
\subtitle{El Futuro y Cómo Seguir Creciendo}
\author{Ricardo Reyes-Jimenez}
\institute{IA para Todos}
\date{\today}



\begin{document}

% Title slide
\begin{frame}
  \titlepage
\end{frame}

% Reset background for normal slides
\usebackgroundtemplate{}



\begin{frame}[fragile]{Objetivo de la Sesión}


\textbf{Al finalizar este módulo, los participantes serán capaces de:}

1. Reconocer que la IA es un lienzo en blanco para cualquier proyecto o afición.
2. Descubrir qué son los agentes o GPTs personalizados.
3. Mantenerse motivados para explorar la IA a su propio ritmo.

\end{frame}



\begin{frame}[fragile]{1. Pasatiempos, Hobbies y Proyectos Personales}


\textbf{No todo es trabajo. La IA es un compañero excelente para tu tiempo libre.}

\begin{itemize}
  \item \textbf{Caso Práctico:} Quieres aprender a tocar la guitarra, escribir una novela, o diseñar un juego de mesa para tus hijos.
  \item \textbf{Uso de IA:} ``Actúa como un profesor de música paciente. Nunca he tocado la guitarra. Diseña un plan de estudio de 4 semanas de 15 minutos diarios, enfocado en aprender a tocar canciones de Los Beatles. Explícame el primer acorde de forma súper sencilla.''
\end{itemize}
\textbf{Resultado:} Aprendizaje personalizado y a tu propio ritmo en cualquier tema.

\end{frame}



\begin{frame}[fragile]{2. GPTs Personalizados y ``Agentes''}


\textbf{Si usas una instrucción todos los días, enséñasela a tu IA de forma permanente.}

\begin{itemize}
  \item \textbf{El Siguiente Nivel:} En lugar de escribir el ``Contexto'' y ``Rol'' cada vez, plataformas como ChatGPT Plus o Claude permiten crear ``Agentes'' (o GPTs) personalizados.
  \item \textbf{Ejemplo Práctico:} Puedes crear un agente llamado ``Revisor de Correos''. Le dices una vez: ``Siempre que te pegue un texto, corrígelo para que suene como un directivo de nivel C, quita las faltas de ortografía, y hazlo breve''.
  \item \textbf{Beneficio:} Desde ese momento, solo pegas el texto y el agente hace el trabajo automáticamente. ¡Tú creas tus propias herramientas!
\end{itemize}
\end{frame}



\begin{frame}[fragile]{3. ¿Hacia dónde vamos?}


\textbf{El verdadero riesgo no es que la IA te quite el trabajo, sino que alguien que use la IA lo haga.}

\begin{itemize}
  \item La tecnología avanza rápido, pero los principios de \textbf{comunicación clara} (Los 5 Pilares, el marco TAM) siempre serán útiles.
  \item \textbf{Tu mayor ventaja competitiva hoy:} La curiosidad. Juega con la herramienta, rompe cosas, haz preguntas tontas y descubre atajos.
  \item \textbf{El límite:} Lo que tú quieras, lo que necesites, y las ganas que tengas de intentarlo.
\end{itemize}
\end{frame}



\begin{frame}[fragile]{Cierre del Curso}


\textbf{Recuerda:}
\begin{itemize}
  \item La IA es una bicicleta para tu mente: te lleva más rápido a donde tú decides ir.
  \item Empieza automatizando algo pequeño hoy (un correo, una lista de compras, un resumen).
  \item Comparte lo que aprendas con otros; así crecemos todos.
\end{itemize}
\textbf{¡Sigan aprendiendo, sigan trabajando y sigan disfrutando de la Inteligencia Artificial!}
\end{frame}




\begin{frame}[allowframebreaks]{References}
  \bibliographystyle{apalike}
  \bibliography{references}
\end{frame}


\end{document}